\documentclass[11pt]{article}
\usepackage[a4paper,margin=27mm]{geometry}
\usepackage{amsmath,amssymb,amsthm,hyperref}
\newtheorem{theorem}{Theorem}
\newtheorem{corollary}{Corollary}
\title{A linear--logarithmic positive-fraction bound from polynomial growth}
\author{Complete proof; independently reviewed}
\date{30 September 2026}
\begin{document}\maketitle
\noindent\textbf{Status.} Full independent mathematical review by the root
agent passed; this is not formal verification.

Let $S$ be a finite set of generators of a group, without assuming
$S=S^{-1}$. Let $A_R$ consist of elements represented by positive words
in $S$ of length at most $R$, including the identity.

\begin{theorem}\label{ore}
Suppose $|A_R|\le(R+1)^D$ for all integers $R\ge0$, where $D\ge1$.
Every signed word of length $L\ge1$ represents a fraction $ab^{-1}$,
where $a,b$ are positive words of lengths at most
\[
 \left\lceil16DL\log(8DL)\right\rceil.
\]
The theorem gives existence; it does not assert an efficient algorithm
for finding the two words.
\end{theorem}
\begin{proof}
Put $a_R=|A_R|$. For every $s\in S$, both $sA_R$ and $A_R$ are
subsets of $A_{R+1}$ and have cardinality $a_R$. Therefore
\[
 |sA_R\mathbin\triangle A_R|\le2(a_{R+1}-a_R).
\]
The same holds for $s^{-1}$ by translation. If $g$ has signed length
$L$, the triangle inequality for symmetric difference, applied to its
successive left translates, gives
\[
 |gA_R\mathbin\triangle A_R|\le2L(a_{R+1}-a_R).
\]
If $a_{R+1}-a_R<a_R/(2L)$, then $gA_R$ intersects $A_R$ (in fact
more than half of each set lies in the intersection). Thus $gb=a$
for some $a,b\in A_R$, giving the required fraction.

It remains to find such an $R$. Set
$T=\lceil16DL\log(8DL)\rceil$. If none exists with $0\le R<T$,
then
\[
 a_T\ge(1+1/(2L))^T\ge\exp(T/(4L)).
\]
On the other hand $a_T\le(T+1)^D$.
For $x=DL\ge1$, the elementary bound
$T+1\le18x\log(8x)\le(8x)^2$ holds; the last inequality follows
from $18\log(8x)\le64x$, whose left-to-right ratio decreases for
$x\ge1$ and is already less than one at $x=1$.
Consequently
\[
 \log a_T\le2D\log(8DL)
 <T/(4L),
\]
a contradiction.
\end{proof}

\begin{corollary}
In the integral reduced free group on $n$ generators, a signed word
of length $L$ has a positive fraction with both lengths
\[
 O\bigl(D_nL\log(D_nL)\bigr),\qquad
 D_n=\sum_{j=1}^n j\frac{n!}{(n-j)!}\le en\,n!.
\]
In particular, signed words of length $\exp(O(n\log n))$ can be
cleared to positive fractions of the same asymptotic length.
\end{corollary}
\begin{proof}
The integral Magnus signature is injective. A positive word of length
at most $R$ has a nonnegative coefficient at each distinct-letter word
of degree $j$, bounded by $R^j$. There are $n!/(n-j)!$ such coordinates.
Therefore the number of possible signatures is at most
\[
 \prod_{j=1}^n(R^j+1)^{n!/(n-j)!}
 \le(R+1)^{D_n}.
\]
Apply Theorem~\ref{ore}. The factorial estimate follows by changing
variables to $n-j$ and bounding the resulting exponential series.
\end{proof}

\noindent\textbf{Effect on the dice problem.} This improves the
positive-fraction stage of the earlier quantitative saturation proof,
which used the Malcev identity in a balanced product tree.
Combined with the subsequently proved
\path{order_doubling_prefix_completion.tex}, it yields an eventual
fair-ray conductor $\exp(O(n\log n))$; see
\path{factorial_order_eventual_conductor.tex}. This is an eventual-size
result and does not provide an $\exp(O(n))$ minimal fair construction.
\end{document}
